\documentclass[12pt, titlepage]{article}

\usepackage{booktabs}
\usepackage{tabularx}
\usepackage{hyperref}
\hypersetup{
    colorlinks,
    citecolor=blue,
    filecolor=black,
    linkcolor=red,
    urlcolor=blue
}
\usepackage[round]{natbib}

\input{../Comments.text}
\input{../Common.text}

\begin{document}

\title{System Verification and Validation Plan for \progname{}} 
\author{\authname}
\date{\today}
	
\maketitle

\pagenumbering{roman}

\section*{Revision History}

\begin{tabularx}{\textwidth}{p{3cm}p{2cm}X}
\toprule {\bf Date} & {\bf Version} & {\bf Notes}\\
\midrule
Date 1 & 1.0 & Notes\\
Date 2 & 1.1 & Notes\\
\bottomrule
\end{tabularx}

~\\
\wss{The intention of the VnV plan is to increase confidence in the software.
However, this does not mean listing every verification and validation technique
that has ever been devised.  The VnV plan should also be a \textbf{feasible}
plan. Execution of the plan should be possible with the time and team available.
If the full plan cannot be completed during the time available, it can either be
modified to ``fake it'', or a better solution is to add a section describing
what work has been completed and what work is still planned for the future.}

\wss{The VnV plan is typically started after the requirements stage, but before
the design stage.  This means that the sections related to unit testing cannot
initially be completed.  The sections will be filled in after the design stage
is complete.  the final version of the VnV plan should have all sections filled
in.}

\newpage

\tableofcontents

\listoftables
\wss{Remove this section if it isn't needed}

\listoffigures
\wss{Remove this section if it isn't needed}

\newpage

\section{Symbols, Abbreviations, and Acronyms}

\renewcommand{\arraystretch}{1.2}
\begin{tabular}{l l} 
  \toprule		
  \textbf{symbol} & \textbf{description}\\
  \midrule 
  T & Test\\
  \bottomrule
\end{tabular}\\

\wss{symbols, abbreviations, or acronyms --- you can simply reference the SRS
  \citep{SRS} tables, if appropriate}

\wss{Remove this section if it isn't needed}

\newpage

\pagenumbering{arabic}

This document ... \wss{provide an introductory blurb and roadmap of the
  Verification and Validation plan}

\section{General Information}

\subsection{Summary}

\wss{Say what software is being tested.  Give its name and a brief overview of
  its general functions.}

\subsection{Objectives}

\wss{State what is intended to be accomplished.  The objective will be around
  the qualities that are most important for your project.  You might have
  something like: ``build confidence in the software correctness,''
  ``demonstrate adequate usability.'' etc.  You won't list all of the qualities,
  just those that are most important.}

\wss{You should also list the objectives that are out of scope.  You don't have 
the resources to do everything, so what will you be leaving out.  For instance, 
if you are not going to verify the quality of usability, state this.  It is also 
worthwhile to justify why the objectives are left out.}

\wss{The objectives are important because they highlight that you are aware of 
limitations in your resources for verification and validation.  You can't do everything, 
so what are you going to prioritize?  As an example, if your system depends on an 
external library, you can explicitly state that you will assume that external library 
has already been verified by its implementation team.}

\subsection{Custom Documents}

\wss{Summarize the CDs (Custom Documents) that were tackled by this project.  CDs 
can include usability testing, code walkthroughs, user documentation, formal
proof, GenderMag personas, Design Thinking, etc.  CDs should have already
been approved by the course instructor as included in your problem statement.
You can use a pull request to update your CDs (in TeamComposition.csv or
Repos.csv) if your plan changes as a result of the VnV planning exercise.}

\subsection{Relevant Documentation}

\wss{Reference relevant documentation.  This will definitely include your SRS
  and your other project documents (design documents, like MG, MIS, etc).  You
  can include these even before they are written, since by the time the project
  is done, they will be written.  You can create BibTeX entries for your
  documents and within those entries include a hyperlink to the documents.}

\citet{SRS}

\wss{Don't just list the other documents.  You should explain why they are relevant and 
how they relate to your VnV efforts.}

\section{Plan}

\wss{Introduce this section.  You can provide a roadmap of the sections to
  come.}

\subsection{Verification and Validation Team}

\wss{Your teammates.  Maybe your supervisor.
  You should do more than list names.  You should say what each person's role is
  for the project's verification.  A table is a good way to summarize this information.}

\subsection{SRS Verification}

\wss{List any approaches you intend to use for SRS verification.  This may
  include ad hoc feedback from reviewers, like your classmates (like your
  primary reviewer), or you may plan for something more rigorous/systematic.}

\wss{If you have a supervisor for the project, you shouldn't just say they will
read over the SRS.  You should explain your structured approach to the review.
Will you have a meeting?  What will you present?  What questions will you ask?
Will you give them instructions for a task-based inspection?  Will you use your
issue tracker?}

\wss{Maybe create an SRS checklist?}

\subsection{Design Verification}

\wss{Plans for design verification}

\wss{The review will include reviews by your classmates}

\wss{Create a checklists?}

\subsection{Verification and Validation Plan Verification}

\wss{The verification and validation plan is an artifact that should also be
verified.  Techniques for this include review and mutation testing.}

\wss{The review will include reviews by your classmates}

\wss{Create a checklists?}

\subsection{Implementation Verification}

\wss{You should at least point to the tests listed in this document and the unit
  testing plan.}

\wss{In this section you would also give any details of any plans for static
  verification of the implementation.  Potential techniques include code
  walkthroughs, code inspection, static analyzers, etc.}

\wss{The final class presentation in CAS 741 could be used as a code
walkthrough.  There is also a possibility of using the final presentation (in
CAS741) for a partial usability survey.}

\subsection{Automated Testing and Verification Tools}

\wss{What tools are you using for automated testing.  Likely a unit testing
  framework and maybe a profiling tool, like ValGrind.  Other possible tools
  include a static analyzer, make, continuous integration tools, test coverage
  tools, etc.  Explain your plans for summarizing code coverage metrics.
  Linters are another important class of tools.  For the programming language
  you select, you should look at the available linters.  There may also be tools
  that verify that coding standards have been respected, like flake9 for
  Python.}

\wss{If you have already done this in the development plan, you can point to
that document.}

\wss{The details of this section will likely evolve as you get closer to the
  implementation.}

\subsection{Software Validation}

\wss{If there is any external data that can be used for validation, you should
  point to it here.  If there are no plans for validation, you should state that
  here.}

\wss{You might want to use review sessions with the stakeholder to check that
the requirements document captures the right requirements.  Maybe task based
inspection?}

\wss{For those capstone teams with an external supervisor, the Rev 0 demo should 
be used as an opportunity to validate the requirements.  You should plan on 
demonstrating your project to your supervisor shortly after the scheduled Rev 0 demo.  
The feedback from your supervisor will be very useful for improving your project.}

\wss{For teams without an external supervisor, user testing can serve the same purpose 
as a Rev 0 demo for the supervisor.}

\wss{This section might reference back to the SRS verification section.}

\section{System Tests}

\wss{There should be text between all headings, even if it is just a roadmap of
the contents of the subsections.}

\subsection{Tests for Functional Requirements}

The tests in this section are grouped into the same six areas as the
functional requirements (FRs) in SRS Section 6: Accounts and Profile, Activity
and Rewards, Simulated Market and Portfolio, Learning and Feedback, Social
Feed, and Reminders and Settings. Each test is derived from the fit criterion
of the FRs it covers, so an FR is verified when all of its tests pass. Every
FR is covered by at least one test, and the traceability table that ends the
System Tests section maps each test back to its FRs.

Each test follows the format of the template, with these conventions:

\begin{itemize}
  \item \textbf{Test ID:} \texttt{test-}\textit{area}\texttt{-}\textit{n},
    where \textit{area} is \texttt{acct}, \texttt{act}, \texttt{mkt},
    \texttt{learn}, \texttt{social}, or \texttt{remind}, matching the SRS
    labels for each area, and \textit{n} numbers the tests within the area.
  \item \textbf{Requirements:} the FRs the test covers, by their SRS number.
  \item \textbf{Control:} \textit{Automatic} tests run with pytest against the
    backend's API, or with Jest for the frontend, on every pull request (see
    Automated Testing and Verification Tools). \textit{Manual} tests are run by
    a tester on an Android or iOS test device, mostly to check what the user
    sees.
\end{itemize}

Automated tests send activity records and market prices to the backend
directly, in place of the health platform and the market data provider, so
every input and expected value is known. The accuracy of those services is out
of scope (see Objectives). Manual tests record activity on a test device
through the phone's health platform, so they also check that \progname{} reads
it correctly.

Tests that depend on the calendar day run in a fixed time zone, with the
backend's clock set to known dates, so that daily totals, streaks, and rest
days can be tested without waiting for real days to pass. Expected values use
the symbolic constants in SRS Appendix C, such as
\texttt{MAX\_DAILY\_CURRENCY}, with their current values given in parentheses.
If a constant changes, the expected outputs of the tests that use it change
with it.

\subsubsection{Activity and Rewards}

These tests cover the activity and rewards requirements in SRS Section 6.2
(FR12 to FR27). Expected currency comes from the daily reward $R_d$ in SRS
Section 6.7, with the values in SRS Appendix C: one unit of currency per 100
steps or wheelchair pushes, one per minute of cycling, swimming, or another
workout, and at most \texttt{MAX\_DAILY\_CURRENCY} (300) per day.

The server calculates rewards only from the daily activity totals the app
sends (Hazard Analysis, SEC-24), and stores no start or end times (SRS Data
Dictionary, Activity Record). Tests of the reward rules therefore run against
the backend with daily totals. Tests that need the time of each activity,
such as assigning it to a calendar day or checking a workout's length, run
against the frontend with Jest. The plausibility limits are input constraints
on activity records, so each limit is tested at its value and just above it.

\begin{enumerate}

\item{\textbf{test-act-1: Earning for each activity type}\\}

Requirements: FR12, FR13

Control: Automatic

Initial State: Six test users, each with a balance of 0, the health notice
acknowledged, no activity goal, and no activity on day $d$.

Input: For day $d$, users 1 to 5 each have a total for one supported type:
2,550 steps, 1,099 wheelchair pushes, 20 minutes of cycling, 15 minutes of
swimming, and 30 minutes of strength training, which counts as another
workout. User 6 has all five totals.

Output: The balances of users 1 to 5 increase by 25, 10, 20, 15, and 30. The
balance of user 6 increases by 100.

Test Case Derivation: Users 1 to 5 each have a single type, so each balance
increase is that type's reward alone, and a type that earns nothing or earns
at the wrong rate shows up on its own (FR12). By $R_d$, each type earns
$\lfloor x_{t,d} / u_t \rfloor$, so 2,550 steps earn 25 and 1,099 pushes earn
10, which also checks that partial units are rounded down. Each minute of a
workout earns 1. User 6 checks that $R_d$ adds up the rewards of all types
(FR13): $25 + 10 + 20 + 15 + 30 = 100$. This is below the daily cap of 300, so
the cap cannot hide a wrong sum.

How test will be performed: A pytest test sends each user's daily totals to
the backend, then reads each balance through the API.

\item{\textbf{test-act-2: Daily earning cap}\\}

Requirements: FR18

Control: Automatic

Initial State: Three test users, each with a balance of 0, the health notice
acknowledged, and no activity on day $d$.

Input: For day $d$, the users have totals of 29,900, 30,000, and 40,000 steps.

Output: The balances increase by 299, 300, and 300.

Test Case Derivation: $R_d = \min(300, \lfloor x_{t,d} / 100 \rfloor)$ gives
299, 300, and $\min(300, 400) = 300$. The three totals check the cap just
below, at, and above \texttt{MAX\_DAILY\_CURRENCY} (300).

How test will be performed: As in test-act-1.

\item{\textbf{test-act-3: Activity counts toward the day in the user's time
zone}\\}

Requirements: FR18

Control: Automatic

Initial State: A test user whose phone is set to the America/Toronto time zone
(UTC$-$5 in January).

Input: Two step records: 20,000 steps ending at 23:30 local time on January
14, 2027 (04:30 UTC on January 15), and 20,000 steps ending at 00:30 local
time on January 15, 2027.

Output: The app's daily totals are 20,000 steps for January 14 and 20,000
steps for January 15.

Test Case Derivation: FR18 counts each activity toward the calendar day in the
user's time zone at the time of the activity. Each day then earns
$\lfloor 20000 / 100 \rfloor = 200$. Grouping the records by UTC day instead
would put 40,000 steps on January 15, which earns only the cap of 300, so the
test detects the wrong grouping.

How test will be performed: A Jest test passes the two records, with their
timestamps and the user's time zone, to the app's code that builds daily
totals, and checks the totals it returns.

\item{\textbf{test-act-4: Day with no activity data}\\}

Requirements: FR16

Control: Automatic

Initial State: A test user with a balance of 100, the health notice
acknowledged, and no activity goal.

Input: The backend receives no activity totals for day $d$, and its clock is
then moved to day $d+1$.

Output: Day $d$ shows a total of 0 for every supported type, and the balance
is still 100.

Test Case Derivation: FR16 treats a day with no data as a day with zero
activity. With $x_{t,d} = 0$ for every type, $R_d = 0$.

How test will be performed: A pytest test reads the totals for day $d$ and the
balance through the API.

\item{\textbf{test-act-5: Daily step and push limits}\\}

Requirements: FR17

Control: Automatic

Initial State: Four test users, each with a balance of 0, the health notice
acknowledged, no activity goal, a streak of 6, and no activity on day $d$.

Input: For day $d$, the users have totals of 150,000 steps, 150,001 steps,
100,000 wheelchair pushes, and 100,001 wheelchair pushes.

Output: The totals at each limit (150,000 steps and 100,000 pushes) are
counted: those users' balances increase by 300, the totals appear in their
activity trends, their streaks become 7, and they reach the 7-day streak
milestone. The totals over each limit earn nothing, do not appear in the
users' trends, reset their streaks to 0, reach no milestone, and are marked as
not counted in the API's response.

Test Case Derivation: Each user has one total, so each limit is checked on its
own. The limits are \texttt{MAX\_DAILY\_\allowbreak STEPS} (150,000) and
\texttt{MAX\_DAILY\_\allowbreak PUSHES} (100,000), and FR17 leaves out a record
only when it exceeds its limit, so a total at the limit is counted and a total
one above it is not. Testing both sides of each limit catches a check that
uses $\geq$ instead of $>$. The totals at the limit earn
more than the cap, so $R_d = 300$. With no goal set, any counted activity
extends the streak (SRS Section 6.7), so a streak of 6 becomes 7. FR25 shows
each milestone in \texttt{MILESTONES} (streaks of 7, 30, and 100 days, and
1,000, 10,000, and 100,000 currency earned in total), so a streak of 7 reaches
the first streak milestone. A total over the limit is left out of
$x_{t,d}$, so $R_d = 0$ and day $d$ does not count toward the streak, which
resets it to 0. Starting from a streak of 6 means that a wrongly counted total
would reach the 7-day milestone, so the test also checks that a total that is
not counted reaches no milestone.

How test will be performed: A pytest test sends each user's totals, then reads
the balance, trends, streak, and each total's status through the API.

\item{\textbf{test-act-6: Workout length limit}\\}

Requirements: FR17

Control: Automatic

Initial State: A test user whose phone has no activity recorded on days $d$
and $d+1$.

Input: A cycling record on day $d$ lasting exactly 12 hours, and a cycling
record on day $d+1$ lasting 12 hours and 1 minute.

Output: The app counts the first workout, giving a cycling total of 720
minutes for day $d$. It leaves the second workout out, giving a cycling total
of 0 for day $d+1$, and marks it as not counted.

Test Case Derivation: \texttt{MAX\_WORKOUT\_DURATION} is 12 hours, and FR17
leaves out only a workout longer than it, so the two records check the limit
at its value and one minute above it. They are on separate days because two
workouts this long cannot fit in one day. Checking a single workout's length
needs its start and end times, which only the phone stores, so the app applies
this limit before it sends the daily total.

How test will be performed: A Jest test passes the records to the app's code
that builds daily totals, and checks the total and each record's status.

\item{\textbf{test-act-7: Showing a record as not counted}\\}

Requirements: FR17

Control: Manual

Initial State: A test user signed in on a test device, with the health notice
acknowledged.

Input: The tester enters a cycling workout of 12 hours and 1 minute for the
previous day in the phone's health platform, then opens \progname{}.

Output: \progname{} shows the workout as not counted, and the user's balance,
activity trends, and streak are the same as before the workout was entered.

Test Case Derivation: FR17 requires the user to be shown that the record was
not counted. test-act-5 and test-act-6 check the underlying values
automatically, and this test checks the whole path on a real device. The
workout is entered for the previous day so that all of it fits within one
finished day, whatever time the test is run.

How test will be performed: The tester notes the balance, trends, and streak,
enters the workout, then checks the same values and the workout's label in
\progname{}.

\item{\textbf{test-act-8: No currency before the health notice is
acknowledged}\\}

Requirements: FR14

Control: Automatic

Initial State: A new test user with a balance of 0 who has not acknowledged
the health notice.

Input: The user has a total of 5,000 steps on day $d$. On day $d+1$, the user
acknowledges the notice, and the app's next sync sends the totals for day $d$
and for day $d+1$, which has 3,000 steps.

Output: The balance is 0 at the end of day $d$ and 80 at the end of day $d+1$.

Test Case Derivation: FR14 allows no currency until the notice is
acknowledged, so day $d$ earns nothing while the notice is unacknowledged.
Each sync re-reads the last \texttt{SYNC\_LOOKBACK\_\allowbreak PERIOD} (7
days) and rewards any record not already rewarded (Hazard Analysis, SEC-20),
so once the notice is acknowledged, day $d$'s 5,000 steps earn 50 alongside
day $d+1$'s $\lfloor 3000 / 100 \rfloor = 30$. The two days have different
totals so that the final balance shows which days were rewarded: 30 would mean
day $d$ was never rewarded, and 80 means both were.

How test will be performed: A pytest test sends day $d$'s total, moves the
backend's clock to day $d+1$, sends the acknowledgement and then both days'
totals, and reads the balance at the end of each day.

\item{\textbf{test-act-9: Health notice shown before the first reward}\\}

Requirements: FR14

Control: Manual

Initial State: A new account signed in on a test device for the first time,
with activity already recorded in the phone's health platform.

Input: The tester opens \progname{}.

Output: Before any currency is awarded, \progname{} shows a notice stating
that it does not give medical advice and that users should consult a health
professional before starting a new exercise program. The balance stays 0 until
the tester acknowledges the notice.

Test Case Derivation: This is the notice and behaviour FR14 requires.
The health platform already has activity when \progname{} is first opened,
so there is currency the app could award straight away, and the test shows
that the notice comes before it.

How test will be performed: The tester compares the notice's text with FR14,
and checks the balance before and after acknowledging it.

\item{\textbf{test-act-10: Rewards without opening the app}\\}

Requirements: FR15

Control: Manual

Initial State: Two test users on separate test devices, each with the health
notice acknowledged and no activity recorded for the current day.

Input: Both testers record a workout of between 10 and 11 minutes with the
phone's own workout tracking. Tester A keeps \progname{} closed for the whole
workout. Tester B opens \progname{} during it. Both open \progname{} after
the workout.

Output: Both balances increase by 10.

Test Case Derivation: At one unit per minute, a workout of between 10 and 11
minutes earns 10, since partial minutes are rounded down. FR15 requires the
same reward whether or not the user interacts with \progname{} during the
activity. The two testers do the workout together, so the only difference
between them is whether they open \progname{} during it.

How test will be performed: Each tester notes their balance before the
workout and compares it with the balance after opening \progname{}.

\item{\textbf{test-act-11: Rewards ignore body weight and calories}\\}

Requirements: FR26

Control: Automatic

Initial State: A test user with a body weight of 70 kg, a balance of 0, the
health notice acknowledged, and no activity on days $d$ to $d+2$.

Input: The user has a total of 30 minutes of cycling on each of days $d$,
$d+1$, and $d+2$. Before day $d+1$, the weight is changed to 90 kg, and before
day $d+2$, it is removed.

Output: The balance increases by 30 on each day, to 90 at the end of day
$d+2$.

Test Case Derivation: The fit criterion of FR26 changes and removes the
user's weight, so this test does the same, with the same activity on each day.
Only the weight differs between the days, so if it affected the reward, at
least one day would earn a different amount. $R_d$ depends only on activity
totals and earning rates, so 30 minutes of cycling earns 30 on each day. The
activity record has no calorie field (SRS Data Dictionary), so calories cannot
reach the reward calculation. Goals are pairs of an activity type and an
amount (SRS Section 6.7), and \texttt{MILESTONES} are streaks of 7, 30, and
100 days and 1,000, 10,000, and 100,000 currency earned in total (SRS Appendix
C), so neither uses weight or calories either.

How test will be performed: A pytest test sends each day's total and the
weight changes, moving the backend's clock between days, and reads the balance
at the end of each day.

\item{\textbf{test-act-12: Goal changes apply from the next day}\\}

Requirements: FR19

Control: Automatic

Initial State: A test user with a daily goal of 10,000 steps, a streak of 3,
the health notice acknowledged, and no activity on days $d$ and $d+1$.

Input: On day $d$, the user lowers the goal to 5,000 steps. The user then has
a total of 6,000 steps on each of days $d$ and $d+1$.

Output: The goal in effect is 10,000 steps on day $d$ and 5,000 steps on day
$d+1$. The streak is 0 at the end of day $d$ and 1 at the end of day $d+1$.

Test Case Derivation: FR19 applies a goal change from the next day, so the
goal in effect on day $d$ is still 10,000. Day $d$'s 6,000 steps miss it, so
day $d$ does not count toward the streak (SRS Section 6.7) and the streak
resets to 0. On day $d+1$, 6,000 steps meet the new goal of 5,000, so the
streak becomes 1. The daily total is between the old and new goals, so
whether each day counts shows which goal was in effect. The goal is lowered,
not raised, because applying changes from the next day is what stops a user
from lowering a goal to keep the current day's streak alive (Hazard Analysis,
SAF-4). The streak starts at 3 so that a change applied on the same day would
give a streak of 4 instead of 0.

How test will be performed: A pytest test sends the goal change and then the
day's total, moves the backend's clock to day $d+1$, sends that day's total,
and reads the goal in effect and the streak for each day through the API.

\item{\textbf{test-act-13: Earning without a goal}\\}

Requirements: FR20

Control: Automatic

Initial State: Two test users, each with a balance of 0, the health notice
acknowledged, and no activity on day $d$. User 1 has no activity goal. User 2
has a daily goal of 3,000 steps.

Input: For day $d$, both users have a total of 4,000 steps.

Output: Both balances increase by 40.

Test Case Derivation: User 1 is the case in the fit criterion of FR20: a user
with no goal who records activity earns currency. User 2 has the same steps
and meets their goal, so comparing the two shows whether a user without a goal
earns less than one who meets a goal. $R_d$ has no goal term (SRS Section
6.7), so both earn $\lfloor 4000 / 100 \rfloor = 40$. A reward that paid extra
for meeting a goal would push users to set goals they did not choose (Hazard
Analysis, SAF-3), and this test would detect one.

How test will be performed: As in test-act-1.

\item{\textbf{test-act-14: No automatic goal changes}\\}

Requirements: FR21

Control: Automatic

Initial State: A test user with a daily goal of 5,000 steps, the health notice
acknowledged, and no activity on days $d$ to $d+6$.

Input: The user has a total of 10,000 steps on each of days $d$ to $d+6$, and
does not change the goal. The backend's clock is then moved to day $d+7$.

Output: The goal in effect on day $d+7$ is still 5,000 steps.

Test Case Derivation: The fit criterion of FR21 has the user exceed their goal
on \texttt{GOAL\_REVIEW\_PERIOD} (7) consecutive days, which are days $d$ to
$d+6$. Each day's total is twice the goal, so the goal is clearly exceeded
throughout, which is when an automatic increase would be most likely. A goal
change applies from the next day (FR19), so a change made at the end of the
period would first take effect on day $d+7$, which is when the goal is read.
The user confirms no change, so any goal other than 5,000 fails the test. A
change the user does confirm is tested in test-act-12.

How test will be performed: A pytest test sends each day's total, moving the
backend's clock between days, and reads the goal in effect on day $d+7$
through the API.

\item{\textbf{test-act-15: Streak with a goal}\\}

Requirements: FR22, FR23

Control: Automatic

Initial State: A test user with a daily goal of 5,000 steps, a streak of 0,
the health notice acknowledged, no rest days, and no activity on days $d$ to
$d+4$.

Input: The user has totals of 6,000, 5,000, and 6,000 steps on days $d$ to
$d+2$. Day $d+3$ is marked as a rest day and has no activity. Day $d+4$ is not
a rest day and has 4,999 steps.

Output: The streak at the end of days $d$ to $d+4$ is 1, 2, 3, 4, and 0.

Test Case Derivation: The days follow the fit criterion of FR22: three days
that meet the goal give a streak of 3, a rest day then gives 4, and a missed
day then gives 0. By SRS Section 6.7, a day counts toward the streak when the
user meets every goal, so $x_{t,d} \geq g$, or when it is a rest day. Day
$d+1$ has exactly the goal and day $d+4$ is one step below it, so the goal is
checked at its value and just below, which catches a check that uses $>$
instead of $\geq$. Day $d+4$ also has activity, which shows that with a goal
set, recording activity is not enough to keep the streak. As in the fit
criterion of FR23, the rest day has no activity, so the streak continues
because of the rest day alone.

How test will be performed: A pytest test sends each day's total or rest day,
moving the backend's clock between days, and reads the streak at the end of
each day through the API.

\item{\textbf{test-act-16: Streak with more than one goal}\\}

Requirements: FR22

Control: Automatic

Initial State: A test user with daily goals of 5,000 steps and 20 minutes of
cycling, a streak of 0, the health notice acknowledged, and no activity on
days $d$ and $d+1$.

Input: The user has 8,000 steps and 20 minutes of cycling on day $d$, and
8,000 steps and 19 minutes of cycling on day $d+1$.

Output: The streak is 1 at the end of day $d$ and 0 at the end of day $d+1$.

Test Case Derivation: SRS Section 6.7 counts a day toward the streak only when
the user meets every goal in effect. Day $d$ meets both goals. On day $d+1$,
the step goal is met but cycling is one minute short, so the streak resets.
The two days differ only in cycling, so the change in the streak comes from
the cycling goal alone. This catches a streak that counts a day when any one
goal is met.

How test will be performed: As in test-act-15.

\item{\textbf{test-act-17: Streak without a goal}\\}

Requirements: FR22

Control: Automatic

Initial State: A test user with no activity goal, a streak of 0, the health
notice acknowledged, and no activity on days $d$ to $d+2$.

Input: The user has 50 steps on day $d$ and 1 minute of swimming on day $d+1$.
The backend receives no activity totals for day $d+2$.

Output: The streak at the end of days $d$ to $d+2$ is 1, 2, and 0.

Test Case Derivation: The fit criterion of FR22 increases the streak of a
user with no goal on each day they record any activity, and SRS Section 6.7
counts such a day when the sum of the day's totals is above 0. The 50 steps
earn no currency, since $\lfloor 50 / 100 \rfloor = 0$, so day $d$ shows that
the streak depends on activity, not on currency earned. Day $d+1$ uses a
different activity type to show that any supported type counts, not only
steps. Day $d+2$ has no activity, which counts as zero (FR16), so the streak
resets.

How test will be performed: As in test-act-15.

\item{\textbf{test-act-18: Rest day limit}\\}

Requirements: FR23

Control: Automatic

Initial State: A test user with no rest days marked, in the week from Monday,
January 4, 2027, to Sunday, January 10, 2027.

Input: The user asks to mark rest days on Tuesday, January 5, Thursday, January
7, Sunday, January 10, and Monday, January 11, 2027. Each request is made on
the day it is for.

Output: The requests for January 5, 7, and 11 are accepted, and the request
for January 10 is refused. The user's rest days are January 5, 7, and 11.

Test Case Derivation: \texttt{REST\_DAYS\_PER\_PERIOD} is 2 and
\texttt{STREAK\_PERIOD} is a calendar week from Monday to Sunday in the user's
time zone (SRS Appendix C). January 5 and 7 use the two rest days of the week
of January 4, so the request for January 10 is a third in the same week and
is refused. January 10 is a Sunday, so the request also checks that Sunday
ends the week that began on Monday. If weeks started on Sunday, it would be
the first rest day of a new week and would be accepted. January 11 is the next
Monday, so it starts a new week and is accepted. A rolling 7-day window would
still contain January 5 and 7 and would refuse it. Each request is made on its
own day so that the test does not depend on whether rest days can be marked
for other days.

How test will be performed: A pytest test moves the backend's clock to each
date, sends the rest day request, checks whether it is accepted, and then
reads the user's rest days through the API.

\end{enumerate}

\subsection{Tests for Nonfunctional Requirements}

\wss{The nonfunctional requirements for accuracy will likely just reference the
  appropriate functional tests from above.  The test cases should mention
  reporting the relative error for these tests.  Not all projects will
  necessarily have nonfunctional requirements related to accuracy.}

\wss{For some nonfunctional tests, you won't be setting a target threshold for
passing the test, but rather describing the experiment you will do to measure
the quality for different inputs.  For instance, you could measure speed versus
the problem size.  The output of the test isn't pass/fail, but rather a summary
table or graph.}

\wss{Tests related to usability could include conducting a usability test and
  survey.  The survey will be in the Appendix.}

\wss{Static tests, review, inspections, and walkthroughs, will not follow the
format for the tests given below.}

\wss{If you introduce static tests in your plan, you need to provide details.
How will they be done?  In cases like code (or document) walkthroughs, who will
be involved? Be specific.}

\subsubsection{Area of Testing1}
		
\paragraph{Title for Test}

\begin{enumerate}

\item{test-id1\\}

Type: Functional, Dynamic, Manual, Static etc.
					
Initial State: 
					
Input/Condition: 
					
Output/Result: 
					
How test will be performed: 
					
\item{test-id2\\}

Type: Functional, Dynamic, Manual, Static etc.
					
Initial State: 
					
Input: 
					
Output: 
					
How test will be performed: 

\end{enumerate}

\subsubsection{Area of Testing2}

...

\subsection{Traceability Between Test Cases and Requirements}

\wss{Provide a table that shows which test cases are supporting which
  requirements.}

\section{Unit Test Description}

\wss{This section should not be filled in until after the MIS (detailed design
  document) has been completed.}

\wss{Reference your MIS (detailed design document) and explain your overall
philosophy for test case selection.}  

\wss{To save space and time, it may be an option to provide less detail in this section.  
For the unit tests you can potentially layout your testing strategy here.  That is, you 
can explain how tests will be selected for each module.  For instance, your test building 
approach could be test cases for each access program, including one test for normal behaviour 
and as many tests as needed for edge cases.  Rather than create the details of the input 
and output here, you could point to the unit testing code.  For this to work, you code 
needs to be well-documented, with meaningful names for all of the tests.}

\subsection{Unit Testing Scope}

\wss{What modules are outside of the scope.  If there are modules that are
  developed by someone else, then you would say here if you aren't planning on
  verifying them.  There may also be modules that are part of your software, but
  have a lower priority for verification than others.  If this is the case,
  explain your rationale for the ranking of module importance.}

\subsection{Tests for Functional Requirements}

\wss{Most of the verification will be through automated unit testing.  If
  appropriate specific modules can be verified by a non-testing based
  technique.  That can also be documented in this section.}

\subsubsection{Module 1}

\wss{Include a blurb here to explain why the subsections below cover the module.
  References to the MIS would be good.  You will want tests from a black box
  perspective and from a white box perspective.  Explain to the reader how the
  tests were selected.}

\begin{enumerate}

\item{test-id1\\}

Type: \wss{Functional, Dynamic, Manual, Automatic, Static etc. Most will
  be automatic}
					
Initial State: 
					
Input: 
					
Output: \wss{The expected result for the given inputs}

Test Case Derivation: \wss{Justify the expected value given in the Output field}

How test will be performed: 
					
\item{test-id2\\}

Type: \wss{Functional, Dynamic, Manual, Automatic, Static etc. Most will
  be automatic}
					
Initial State: 
					
Input: 
					
Output: \wss{The expected result for the given inputs}

Test Case Derivation: \wss{Justify the expected value given in the Output field}

How test will be performed: 

\item{...\\}
    
\end{enumerate}

\subsubsection{Module 2}

...

\subsection{Tests for Nonfunctional Requirements}

\wss{If there is a module that needs to be independently assessed for
  performance, those test cases can go here.  In some projects, planning for
  nonfunctional tests of units will not be that relevant.}

\wss{These tests may involve collecting performance data from previously
  mentioned functional tests.}

\subsubsection{Module ?}
		
\begin{enumerate}

\item{test-id1\\}

Type: \wss{Functional, Dynamic, Manual, Automatic, Static etc. Most will
  be automatic}
					
Initial State: 
					
Input/Condition: 
					
Output/Result: 
					
How test will be performed: 
					
\item{test-id2\\}

Type: Functional, Dynamic, Manual, Static etc.
					
Initial State: 
					
Input: 
					
Output: 
					
How test will be performed: 

\end{enumerate}

\subsubsection{Module ?}

...

\subsection{Traceability Between Test Cases and Modules}

\wss{Provide evidence that all modules have been considered.}
				
\bibliographystyle{plainnat}

\bibliography{../../refs/References}

\newpage

\section{Appendix}

This is where you can place additional information.

\subsection{Symbolic Parameters}

The definition of the test cases will call for SYMBOLIC\_CONSTANTS.
Their values are defined in this section for easy maintenance.

\subsection{Usability Survey Questions?}

\wss{This is a section that would be appropriate for some projects.}

\newpage{}

\section{Appendix --- Generative AI Use Disclosure}

\wss{ChatGPT (version 5) was used to brainstorm ideas for the template for this disclose section.}

\subsection{AI Tools Used}

\wss{Give the name of the AI tool(s) used and the version.}

\subsection{How and Where Used}

\wss{Explain what the tool(s) were used for.  Examples include: brainstorming,
explaining a technical concept, reviewing the document, generating or modifying
text, generating or modifying diagrams, identifying errors or omissions or
inconsistencies.}

\wss{For each of the uses, identify which parts of the document were affected.}

\wss{You do not need to report every interaction, but you should disclose the 
uses that materially contributed to your work}

\subsection{Verification of AI-Generated Content}

\wss{\begin{itemize}
    \item Describe how the team checked AI-generated or AI-assisted material
    before including it in the document.
    \item In particular, verify factual claims, technical assertions,
    requirements, calculations, references, and other information for which an
    AI system may produce incorrect or fabricated results.
    \item \textbf{The team is responsible for the accuracy, completeness,
consistency, and appropriateness of the submitted document, regardless of
whether material was generated by a human or an AI system.}
\end{itemize}}

\newpage{}

\section{Appendix --- Reflection}

\wss{This section is not required for CAS 741}

The information in this section will be used to evaluate the team members on the
graduate attribute of Lifelong Learning.

\input{../Reflection.text}

\begin{enumerate}
  \item What went well while writing this deliverable? 
  \item What pain points did you experience during this deliverable, and how
    did you resolve them?
  \item What knowledge and skills will the team collectively need to acquire to
  successfully complete the verification and validation of your project?
  Examples of possible knowledge and skills include dynamic testing knowledge,
  static testing knowledge, specific tool usage, Valgrind etc.  You should look to
  identify at least one item for each team member.
  \item For each of the knowledge areas and skills identified in the previous
  question, what are at least two approaches to acquiring the knowledge or
  mastering the skill?  Of the identified approaches, which will each team
  member pursue, and why did they make this choice?
\end{enumerate}

\end{document}